\makeatletter
\def\input@path{{./}{../}{../../}{../../../}{../../../../}{preamble/}{../preamble/}{../../preamble/}{../../../preamble/}{../../../../preamble/}}
\makeatother
\input{preamble_exams.tex}
\graphicspath{{./}{../}{../../}{../../../}{../../../../}{preamble/}{../preamble/}{../../preamble/}{../../../preamble/}{../../../../preamble/}}
\linespread{1.0}
\setlength{\parskip}{5pt}
\usepackage{needspace}
\title{Grade 11 --- Learning Evidence 3\\C8, C9 and C10 --- Exam Solutions --- Version 3\\\large 1st Term \quad 2026--2027}
\author{}
\date{}
\begin{document}
\maketitle
\setcounter{tocdepth}{1}
\setcounter{secnumdepth}{2}
\phantomsection
\label{toc}
\tableofcontents

\section*{Learning objective}
Apply continuous probability models and sampling distributions to investigate economic and financial questions, justify statistical methods, and communicate evidence-based conclusions.

\section*{Criteria assessment}
Partial credit may be awarded when the correct procedure is clearly demonstrated, even if arithmetic errors are made. No points will be awarded for answers without a demonstrated procedure.

A minimum of 4 points is required to pass each criterion. A 1-point penalty will be applied to each criterion if the student's name is not written on the assessment or if the assessment is submitted late.


\begin{center}
\begin{tabular}{|c|c|c|c|c|}
\hline
\textbf{Questions} & \textbf{C8} & \textbf{C9} & \textbf{C10} & \textbf{Total}\\\hline
1--10 & 10 & -- & -- & 10\\\hline
11--16 & -- & 12 & -- & 12\\\hline
17 & -- & -- & 6 & 6\\\hline
\textbf{Total} & \textbf{10} & \textbf{12} & \textbf{6} & \textbf{28}\\\hline
\end{tabular}
\end{center}


The worked solutions show the reasoning using the supplied normal and Student's $t$ reference sheets; no interpolation, CDF calculator, or inverse-CDF calculator is needed.
C8 is worth \textbf{10 points}, C9 \textbf{12 points}, and C10 \textbf{6 points}: \textbf{28 points total}.
At least \textbf{4 points per criterion} are required to pass that criterion.
All calculations within one problem share its stated allocation, including both interval endpoints.

\section*{Evaluated criteria}
\begin{enumerate}[label=\textbf{C\arabic*},start=8,leftmargin=3em]
\item Uses Student's t-distribution tables, degrees of freedom, and one-tailed or two-tailed areas to determine critical values, central probabilities, and tail probabilities or probability ranges.
\item Calculates and compares critical values and probability areas from the standard normal and Student's t-distributions, selecting and justifying the appropriate distribution according to whether the population standard deviation is known or estimated from the sample.
\item Integrates continuous probability modelling, simulation, sampling analysis, and portfolio-risk measures to investigate an economic or financial question, justifying the selected methods and communicating results, assumptions, limitations, and evidence-based conclusions.
\end{enumerate}

\section*{Models and table conventions}
Problems 1--10 use the supplied signed-deviation model $X=m+aT$, $T\sim t_\nu$.
Problems 11--16 compare
\[
 X_T=m+aT,\quad T\sim t_\nu,\qquad X_N=m+aZ,\quad Z\sim N(0,1),
 \qquad \boxed{\nu=n-1}.
\]
Here $n$ is the number of sample observations, $m$ is the location (and mean for $\nu>1$), and $a>0$ is the scale.
Use the stated \emph{approximate standard deviation} directly as $a$, without adjustment.
This is a teaching approximation:
\[
 \operatorname{SD}(X_T)=a\sqrt{\frac{\nu}{\nu-2}}\quad(\nu>2),
 \qquad \operatorname{SD}(X_N)=a.
\]
The two models share a scale, not an actual variance. The approximation is coarser for small $\nu$ and improves as $\nu$ increases.
Negative deviations mean below the contextual reference (early or under budget); positive deviations mean above it (late or over budget).
Inverse intervals enclose modelled deviations. Standardize and back-transform using
\[
 \boxed{c=\frac{x-m}{a}},\qquad \boxed{x=m+ac}.
\]
The year-local sheets are in \texttt{book/fundamentals/normal-table.tex}.
In the $t$ table, row $\nu$ and one-tail column $\alpha$ give positive $t^*$ satisfying
\[
 P(T>t^*)=\alpha,\qquad P(T<-t^*)=\alpha,\qquad
 P(-t^*<T<t^*)=1-2\alpha.
\]
Obtain negative cutoffs by symmetry. The printed columns are
$0.30,0.20,0.10,0.05,0.025,0.01,0.005$.
Direct bounds recover printed three-decimal entries exactly through standardization.
Report C8 probabilities as implied by these entries, and contextual cutoffs to two decimals.

The normal table gives four-decimal cumulative areas $\Phi(z)=P(Z\le z)$ for positive $z$.
Use $\Phi(-z)=1-\Phi(z)$. For direct normal probabilities, round each standardized boundary to two decimals \emph{before} lookup.
For inverse normal problems, select the closest printed CDF entry.
At targets $0.0500$ and $0.9500$, magnitudes $1.64$ and $1.65$ tie; choose $|z|=1.64$, as stipulated for this exam.
Use the CDF table for inverse normal lookups.

For \textbf{every C9 problem}, calculate both results and
\[
 E_A=|r_T-r_N|,\qquad E_{Pt}=\frac{100E_A}{|r_T|},\qquad
 E_{Pn}=\frac{100E_A}{|r_N|}.
\]
For direct problems, $r_T,r_N$ are probabilities; for inverse problems, they are contextual cutoffs after back-transformation.
Compare inverse interval endpoints separately.
A zero denominator makes that percentage undefined; the data below avoid zero denominators.
These measures compare model results and do not establish error against observed truth.
Report C9 probabilities and their $E_A$ to four decimals; contextual cutoffs and their $E_A$ to two decimals with units; and percentage comparisons to two decimals.
Use unrounded back-transformed cutoffs for inverse comparisons, rounding only final results.

All results are table-based approximations, not exact distribution values: printed critical values and CDF entries are rounded.
Equalities in the worked arithmetic use selected table values as working values.
Both models are continuous, so endpoints have probability zero.
The formally stated criteria above are retained; the supplied activity models govern the C8/C9 calculations here.


\needspace{24\baselineskip}
\sectionScored{Retail revenue alerts}{C8-\sml{1}}{sec:q1}
{
\textbf{C8 -- Direct unequal combined tails.} Over $n=10$ weeks, revenue deviation from plan, in thousands of dollars, has mean $m=2$ and approximate standard deviation $a=4$. Use $X=2+4T$, $T\sim t_{n-1}$.
\QuestionsTasks[sec:q1]{
{8}/{sol:q1}/{Find the probability of an alert: $X<-7.048$ or $X>7.532$. \score{1}}
}
\needspace{5\baselineskip}
\phantomsection
\label{sol:q1}
\subsection*{C8 -- Problem 1: Direct unequal combined tails solution \score{1}}
Here $\nu=10-1=9$. Standardize:
\[
\begin{aligned}
c_L&=\frac{-7.048-(2)}{4}=-2.262.\\
c_U&=\frac{7.532-(2)}{4}=1.383.
\end{aligned}
\]
Row $9$, column $\alpha=0.025$, gives $t^*=2.262$. Row $9$, column $\alpha=0.1$, gives $t^*=1.383$. 
By symmetry, the left tail has the first printed area. Add the disjoint tails:
\[
P(X<-7.048\text{ or }X>7.532)\approx0.025+0.1=\boxed{0.125}.
\]
About $12.5\%$ of weeks trigger an alert: revenue is more than $7.05$ thousand dollars below plan or $7.53$ thousand dollars above it.
\secInicio[][sec:q1]
}

\needspace{24\baselineskip}
\sectionScored{Warehouse energy deviations}{C8-\sml{1}}{sec:q2}
{
\textbf{C8 -- Direct interval above the mean.} Over $n=21$ days, electricity-use deviation from forecast, in kWh, has mean $m=3$ and approximate standard deviation $a=6$. Use $X=3+6T$, $T\sim t_{n-1}$.
\QuestionsTasks[sec:q2]{
{8}/{sol:q2}/{Find $P(6.198<X<10.95)$. \score{1}}
}
\needspace{5\baselineskip}
\phantomsection
\label{sol:q2}
\subsection*{C8 -- Problem 2: Direct interval above the mean solution \score{1}}
Here $\nu=21-1=20$. Standardize:
\[
\begin{aligned}
c_L&=\frac{6.198-(3)}{6}=0.533.\\
c_U&=\frac{10.95-(3)}{6}=1.325.
\end{aligned}
\]
Row $20$, column $\alpha=0.3$, gives $t^*=0.533$. Row $20$, column $\alpha=0.1$, gives $t^*=1.325$. 
Both bounds exceed $m=3$. Subtract the nested right tails:
\[
P(6.198<X<10.95)\approx0.3-0.1=\boxed{0.2}.
\]
About $20\%$ of days use between $6.20$ and $10.95$ kWh above forecast.
\secInicio[][sec:q2]
}

\needspace{24\baselineskip}
\sectionScored{Courier arrival deviations}{C8-\sml{1}}{sec:q3}
{
\textbf{C8 -- Direct left tail.} A sample of $n=13$ courier trips has arrival-time deviation from schedule, in minutes, with mean $m=0.8$ and approximate standard deviation $a=4$. Use $X=0.8+4T$, $T\sim t_{n-1}$.
\QuestionsTasks[sec:q3]{
{8}/{sol:q3}/{Find $P(X<-4.624)$. \score{1}}
}
\needspace{5\baselineskip}
\phantomsection
\label{sol:q3}
\subsection*{C8 -- Problem 3: Direct left tail solution \score{1}}
Here $\nu=13-1=12$. Standardize:
\[
\begin{aligned}
c&=\frac{-4.624-(0.8)}{4}=-1.356.
\end{aligned}
\]
Row $12$, column $\alpha=0.1$, gives $t^*=1.356$. 
By symmetry,
\[
P(T<-1.356)\approx\boxed{0.1}.
\]
About $10\%$ of trips arrive more than $4.62$ minutes early.
\secInicio[][sec:q3]
}

\needspace{24\baselineskip}
\sectionScored{Freight cost deviations}{C8-\sml{1}}{sec:q4}
{
\textbf{C8 -- Direct right tail.} A sample of $n=19$ consignments has cost deviation from budget, in dollars, with mean $m=-4$ and approximate standard deviation $a=8$. Use $X=-4+8T$, $T\sim t_{n-1}$.
\QuestionsTasks[sec:q4]{
{8}/{sol:q4}/{Find $P(X>-8.272)$. \score{1}}
}
\needspace{5\baselineskip}
\phantomsection
\label{sol:q4}
\subsection*{C8 -- Problem 4: Direct right tail solution \score{1}}
Here $\nu=19-1=18$. Standardize:
\[
\begin{aligned}
c&=\frac{-8.272-(-4)}{8}=-0.534.
\end{aligned}
\]
Row $18$, column $\alpha=0.3$, gives $t^*=0.534$. 
Symmetry gives the excluded left tail; take its complement:
\[
1-P(T<-0.534)\approx1-0.3=\boxed{0.7}.
\]
About $70\%$ have costs above the threshold of approximately $\$8.27$ below budget. The cutoff is below $m=-4$, so the right-tail event includes more than half the symmetric model.
\secInicio[][sec:q4]
}

\needspace{24\baselineskip}
\sectionScored{Crop shipment deviations}{C8-\sml{1}}{sec:q5}
{
\textbf{C8 -- Direct interval below the mean.} A sample of $n=12$ crop shipments has mass deviation from forecast, in tonnes, with mean $m=-0.8$ and approximate standard deviation $a=3$. Use $X=-0.8+3T$, $T\sim t_{n-1}$.
\QuestionsTasks[sec:q5]{
{8}/{sol:q5}/{Find $P(-6.188<X<-2.42)$. \score{1}}
}
\needspace{5\baselineskip}
\phantomsection
\label{sol:q5}
\subsection*{C8 -- Problem 5: Direct interval below the mean solution \score{1}}
Here $\nu=12-1=11$. Standardize:
\[
\begin{aligned}
c_L&=\frac{-6.188-(-0.8)}{3}=-1.796.\\
c_U&=\frac{-2.42-(-0.8)}{3}=-0.540.
\end{aligned}
\]
Row $11$, column $\alpha=0.05$, gives $t^*=1.796$. Row $11$, column $\alpha=0.3$, gives $t^*=0.540$. 
Both bounds lie below $m=-0.8$. By symmetry the cumulative areas are $0.05$ and $0.3$:
\[
P(-6.188<X<-2.42)\approx0.3-0.05=\boxed{0.25}.
\]
About $25\%$ of shipments fall short of forecast by between $2.42$ and $6.19$ tonnes.
\secInicio[][sec:q5]
}

\needspace{24\baselineskip}
\sectionScored{Operating cost cutoff}{C8-\sml{1}}{sec:q6}
{
\textbf{C8 -- Inverse lower-tail cutoff.} Over $n=20$ production days, operating-cost deviation from budget, in dollars, has mean $m=-4.2$ and approximate standard deviation $a=5.5$. Use $X=-4.2+5.5T$, $T\sim t_{n-1}$.
\QuestionsTasks[sec:q6]{
{8}/{sol:q6}/{Find $x$ such that $P(X<x)=0.10$. \score{1}}
}
\needspace{5\baselineskip}
\phantomsection
\label{sol:q6}
\subsection*{C8 -- Problem 6: Inverse lower-tail cutoff solution \score{1}}
Here $\nu=20-1=19$. Standardize the unknown cutoff(s) using $c=(x-(-4.2))/5.5$. The lower-tail target is $0.1$.
Row $19$, column $\alpha=0.1$, gives $t^*=1.328$. 
Obtain the negative cutoff by symmetry, giving $c=-1.328$.
Back-transform, retaining unrounded values:
\[
\begin{aligned}
x&=-4.2+5.5(-1.328)=-11.504\approx\boxed{-11.50\text{ dollars}}.
\end{aligned}
\]
About $10\%$ of days have costs more than $\$11.50$ below budget.
\secInicio[][sec:q6]
}

\needspace{24\baselineskip}
\sectionScored{Package mass alerts}{C8-\sml{1}}{sec:q7}
{
\textbf{C8 -- Inverse equal-tail alert limits.} A sample of $n=15$ packages has mass deviation from target, in grams, with mean $m=-0.7$ and approximate standard deviation $a=1.8$. Use $X=-0.7+1.8T$, $T\sim t_{n-1}$.
\QuestionsTasks[sec:q7]{
{8}/{sol:q7}/{Find two alert limits leaving a combined outside probability of $0.05$, split equally. \score{1}}
}
\needspace{5\baselineskip}
\phantomsection
\label{sol:q7}
\subsection*{C8 -- Problem 7: Inverse equal-tail alert limits solution \score{1}}
Here $\nu=15-1=14$. Standardize the unknown cutoff(s) using $c=(x-(-0.7))/1.8$. Each tail has probability $0.05/2=0.025$.
Row $14$, column $\alpha=0.025$, gives $t^*=2.145$. 
Obtain the negative cutoff by symmetry, giving $c=-2.145$, $c=2.145$.
Back-transform, retaining unrounded values:
\[
\begin{aligned}
L&=-0.7+1.8(-2.145)=-4.561\approx\boxed{-4.56\text{ g}}.\\
U&=-0.7+1.8(2.145)=3.161\approx\boxed{3.16\text{ g}}.
\end{aligned}
\]
About $2.5\%$ of packages fall below $L$ and $2.5\%$ above $U$. An alert means more than $4.56$ g under target or $3.16$ g over target.
\secInicio[][sec:q7]
}

\needspace{24\baselineskip}
\sectionScored{Dispatch time alerts}{C8-\sml{1}}{sec:q8}
{
\textbf{C8 -- Inverse unequal-tail alert limits.} A sample of $n=23$ dispatches has departure-time deviation from schedule, in minutes, with mean $m=-1.8$ and approximate standard deviation $a=6$. Use $X=-1.8+6T$, $T\sim t_{n-1}$.
\QuestionsTasks[sec:q8]{
{8}/{sol:q8}/{Find $L,U$ with $P(X<L)=0.01$ and $P(X>U)=0.10$. \score{1}}
}
\needspace{5\baselineskip}
\phantomsection
\label{sol:q8}
\subsection*{C8 -- Problem 8: Inverse unequal-tail alert limits solution \score{1}}
Here $\nu=23-1=22$. Standardize the unknown cutoff(s) using $c=(x-(-1.8))/6$. Use the two unequal tails separately.
Row $22$, column $\alpha=0.01$, gives $t^*=2.508$. Row $22$, column $\alpha=0.1$, gives $t^*=1.321$. 
Obtain the negative cutoff by symmetry, giving $c=-2.508$, $c=1.321$.
Back-transform, retaining unrounded values:
\[
\begin{aligned}
L&=-1.8+6(-2.508)=-16.848\approx\boxed{-16.85\text{ minutes}}.\\
U&=-1.8+6(1.321)=6.126\approx\boxed{6.13\text{ minutes}}.
\end{aligned}
\]
About $1\%$ depart more than $16.85$ minutes early, and $10\%$ more than $6.13$ minutes late. The combined alert probability is $0.01+0.10=0.11$; $0.89$ lies between the limits.
\secInicio[][sec:q8]
}

\needspace{24\baselineskip}
\sectionScored{Loading time cutoff}{C8-\sml{1}}{sec:q9}
{
\textbf{C8 -- Inverse upper-tail cutoff.} A sample of $n=26$ loading jobs has duration deviation from target, in minutes, with mean $m=3$ and approximate standard deviation $a=7$. Use $X=3+7T$, $T\sim t_{n-1}$.
\QuestionsTasks[sec:q9]{
{8}/{sol:q9}/{Find $x$ such that $P(X>x)=0.025$. \score{1}}
}
\needspace{5\baselineskip}
\phantomsection
\label{sol:q9}
\subsection*{C8 -- Problem 9: Inverse upper-tail cutoff solution \score{1}}
Here $\nu=26-1=25$. Standardize the unknown cutoff(s) using $c=(x-(3))/7$. The right-tail target is $0.025$.
Row $25$, column $\alpha=0.025$, gives $t^*=2.060$. 
The positive cutoff is $c=2.060$.
Back-transform, retaining unrounded values:
\[
\begin{aligned}
x&=3+7(2.060)=17.42\approx\boxed{17.42\text{ minutes}}.
\end{aligned}
\]
About $2.5\%$ of jobs take more than $17.42$ minutes longer than target.
\secInicio[][sec:q9]
}

\needspace{24\baselineskip}
\sectionScored{Demand deviation interval}{C8-\sml{1}}{sec:q10}
{
\textbf{C8 -- Inverse central interval.} Over $n=17$ weeks, demand deviation from forecast, in tonnes, has mean $m=4$ and approximate standard deviation $a=9$. Use $X=4+9T$, $T\sim t_{n-1}$.
\QuestionsTasks[sec:q10]{
{8}/{sol:q10}/{Find $L<U$ enclosing the central $90\%$ of deviations with equal excluded tails. \score{1}}
}
\needspace{5\baselineskip}
\phantomsection
\label{sol:q10}
\subsection*{C8 -- Problem 10: Inverse central interval solution \score{1}}
Here $\nu=17-1=16$. Standardize the unknown cutoff(s) using $c=(x-(4))/9$. Each tail is $(1-0.90)/2=0.05$.
Row $16$, column $\alpha=0.05$, gives $t^*=1.746$. 
Obtain the negative cutoff by symmetry, giving $c=-1.746$, $c=1.746$.
Back-transform, retaining unrounded values:
\[
\begin{aligned}
L&=4+9(-1.746)=-11.714\approx\boxed{-11.71\text{ tonnes}}.\\
U&=4+9(1.746)=19.714\approx\boxed{19.71\text{ tonnes}}.
\end{aligned}
\]
About $90\%$ lie between $11.71$ tonnes below forecast and $19.71$ tonnes above it, with $5\%$ in each excluded tail.
\secInicio[][sec:q10]
}

\needspace{38\baselineskip}
\sectionScored{Sales revenue shortfalls}{C9-\sml{2}}{sec:q11}
{
\textbf{C9 -- Direct left tail.} Over $n=17$ weeks, sales-revenue deviation from target, in thousands of dollars, has mean $m=-1.8$ and approximate standard deviation $a=5$. Compare $X_T=-1.8+5T$, $T\sim t_{n-1}$, with $X_N=-1.8+5Z$.
\QuestionsTasks[sec:q11]{
{9}/{sol:q11}/{Find $P(X<-12.4)$. Find both model results, calculate $E_A$, $E_{Pt}$, and $E_{Pn}$. \score{2}}
}
\needspace{5\baselineskip}
\phantomsection
\label{sol:q11}
\subsection*{C9 -- Problem 11: Direct left tail solution \score{2}}
Here $\nu=17-1=16$. Standardize:
\[
\begin{aligned}
c&=\frac{-12.4-(-1.8)}{5}=-2.120.
\end{aligned}
\]
Row $16$, column $\alpha=0.025$, gives $t^*=2.120$. 
By symmetry,
\[
p_T=P(T<-2.120)\approx\boxed{0.0250}.
\]
For normal, round the standardized bound(s) to $z=-2.12$.
Row $2.1$, column $0.02$, gives $\Phi(2.12)=0.9830$. 
\[
p_N=1-0.9830=\boxed{0.0170}.
\]
\[
\begin{aligned}
E_{A}&=|0.025-(0.017)|=0.008\approx\boxed{0.0080},\\
E_{Pt}&=\frac{100(0.008)}{|0.025|}\approx\boxed{32.00\%},\\
E_{Pn}&=\frac{100(0.008)}{|0.017|}\approx\boxed{47.06\%}.
\end{aligned}
\]
A shortfall greater than $12.40$ thousand dollars has probability about $2.50\%$ under $t$, compared with $1.70\%$ under normal, a difference of $0.80$ percentage points.
\secInicio[][sec:q11]
}

\needspace{38\baselineskip}
\sectionScored{Fuel forecast deviations}{C9-\sml{2}}{sec:q12}
{
\textbf{C9 -- Direct interval crossing the mean.} Over $n=19$ operating days, fuel-use deviation from forecast, in litres, has mean $m=0.9$ and approximate standard deviation $a=6$. Compare $X_T=0.9+6T$, $T\sim t_{n-1}$, with $X_N=0.9+6Z$.
\QuestionsTasks[sec:q12]{
{9}/{sol:q12}/{Find $P(-4.272<X<11.304)$. Find both model results, calculate $E_A$, $E_{Pt}$, and $E_{Pn}$. \score{2}}
}
\needspace{5\baselineskip}
\phantomsection
\label{sol:q12}
\subsection*{C9 -- Problem 12: Direct interval crossing the mean solution \score{2}}
Here $\nu=19-1=18$. Standardize:
\[
\begin{aligned}
c_L&=\frac{-4.272-(0.9)}{6}=-0.862.\\
c_U&=\frac{11.304-(0.9)}{6}=1.734.
\end{aligned}
\]
Row $18$, column $\alpha=0.2$, gives $t^*=0.862$. Row $18$, column $\alpha=0.05$, gives $t^*=1.734$. 
The interval crosses $m=0.9$ and excludes unequal tails:
\[
p_T\approx1-0.2-0.05=\boxed{0.7500}.
\]
For normal, round the standardized bound(s) to $z=-0.86$, $z=1.73$.
Row $0.8$, column $0.06$, gives $\Phi(0.86)=0.8051$. Row $1.7$, column $0.03$, gives $\Phi(1.73)=0.9582$. 
\[
p_N=0.9582-(1-0.8051)=\boxed{0.7633}.
\]
\[
\begin{aligned}
E_{A}&=|0.75-(0.7633)|=0.0133\approx\boxed{0.0133},\\
E_{Pt}&=\frac{100(0.0133)}{|0.75|}\approx\boxed{1.77\%},\\
E_{Pn}&=\frac{100(0.0133)}{|0.7633|}\approx\boxed{1.74\%}.
\end{aligned}
\]
Fuel use lies from $4.27$ litres below forecast to $11.30$ litres above it with probability about $75.00\%$ under $t$ or $76.33\%$ under normal. Normal assigns $1.33$ percentage points more to this interval.
\secInicio[][sec:q12]
}

\needspace{38\baselineskip}
\sectionScored{Delivery schedule deviations}{C9-\sml{2}}{sec:q13}
{
\textbf{C9 -- Direct right tail.} A sample of $n=16$ deliveries has arrival-time deviation from schedule, in minutes, with mean $m=2.6$ and approximate standard deviation $a=5$. Compare $X_T=2.6+5T$, $T\sim t_{n-1}$, with $X_N=2.6+5Z$.
\QuestionsTasks[sec:q13]{
{9}/{sol:q13}/{Find $P(X>-0.08)$. Find both model results, calculate $E_A$, $E_{Pt}$, and $E_{Pn}$. \score{2}}
}
\needspace{5\baselineskip}
\phantomsection
\label{sol:q13}
\subsection*{C9 -- Problem 13: Direct right tail solution \score{2}}
Here $\nu=16-1=15$. Standardize:
\[
\begin{aligned}
c&=\frac{-0.08-(2.6)}{5}=-0.536.
\end{aligned}
\]
Row $15$, column $\alpha=0.3$, gives $t^*=0.536$. 
Symmetry gives the excluded left tail; take its complement:
\[
p_T=1-P(T<-0.536)\approx1-0.3=\boxed{0.7000}.
\]
For normal, round the standardized bound(s) to $z=-0.54$.
Row $0.5$, column $0.04$, gives $\Phi(0.54)=0.7054$. 
\[
\Phi(-0.54)=1-0.7054,\quad p_N=1-(1-0.7054)=\boxed{0.7054}.
\]
\[
\begin{aligned}
E_{A}&=|0.7-(0.7054)|=0.0054\approx\boxed{0.0054},\\
E_{Pt}&=\frac{100(0.0054)}{|0.7|}\approx\boxed{0.77\%},\\
E_{Pn}&=\frac{100(0.0054)}{|0.7054|}\approx\boxed{0.77\%}.
\end{aligned}
\]
The event includes arrivals less than about $0.08$ minutes early and all late arrivals. Normal assigns $70.54\%$, which is $0.54$ percentage points higher than $t$'s $70.00\%$. Both exceed $0.5$ because the cutoff is below $m=2.6$; $t$ does not give the larger probability for this event.
\secInicio[][sec:q13]
}

\needspace{49\baselineskip}
\sectionScored{Cargo mass deviation interval}{C9-\sml{2}}{sec:q14}
{
\textbf{C9 -- Inverse central interval.} A sample of $n=14$ cargo loads has mass deviation from target, in kg, with mean $m=2.2$ and approximate standard deviation $a=4.5$. Compare $X_T=2.2+4.5T$, $T\sim t_{n-1}$, with $X_N=2.2+4.5Z$.
\QuestionsTasks[sec:q14]{
{9}/{sol:q14}/{Find the central $95\%$ interval with equal excluded tails under each model. Calculate $E_A$, $E_{Pt}$, and $E_{Pn}$ separately for each endpoint. \score{2}}
}
\needspace{5\baselineskip}
\phantomsection
\label{sol:q14}
\subsection*{C9 -- Problem 14: Inverse central interval solution \score{2}}
Here $\nu=14-1=13$. Standardize the unknown cutoff(s) using $c=(x-(2.2))/4.5$. Each tail is $(1-0.95)/2=0.025$, so cumulative targets are $0.0250$ and $0.9750$.
Row $13$, column $\alpha=0.025$, gives $t^*=2.160$. 
Obtain the negative cutoff by symmetry, giving $c=-2.160$, $c=2.160$.
Row $1.9$, column $0.06$, gives $\Phi(1.96)=0.9750$. This matches $0.9750$ at printed precision; symmetry gives $\Phi(-1.96)=0.0250$.
Back-transform, retaining unrounded values:
\[
\begin{aligned}
L_T&=2.2+4.5(-2.160)=-7.52\approx\boxed{-7.52\text{ kg}}.\\
U_T&=2.2+4.5(2.160)=11.92\approx\boxed{11.92\text{ kg}}.
\end{aligned}
\]
\[
\begin{aligned}
L_N&=2.2+4.5(-1.96)=-6.62\approx\boxed{-6.62\text{ kg}}.\\
U_N&=2.2+4.5(1.96)=11.02\approx\boxed{11.02\text{ kg}}.
\end{aligned}
\]
\textbf{Lower endpoint comparison:}
\[
\begin{aligned}
E_{A,L}&=|-7.52-(-6.62)|=0.9\approx\boxed{0.90\text{ kg}},\\
E_{Pt,L}&=\frac{100(0.9)}{|-7.52|}\approx\boxed{11.97\%},\\
E_{Pn,L}&=\frac{100(0.9)}{|-6.62|}\approx\boxed{13.60\%}.
\end{aligned}
\]
\textbf{Upper endpoint comparison:}
\[
\begin{aligned}
E_{A,U}&=|11.92-(11.02)|=0.9\approx\boxed{0.90\text{ kg}},\\
E_{Pt,U}&=\frac{100(0.9)}{|11.92|}\approx\boxed{7.55\%},\\
E_{Pn,U}&=\frac{100(0.9)}{|11.02|}\approx\boxed{8.17\%}.
\end{aligned}
\]
The central intervals are approximately $(-7.52,11.92)$ kg under $t$ and $(-6.62,11.02)$ kg under normal, with $2.5\%$ in each excluded tail. Negative limits mean below target mass; positive limits mean above it. The normal interval extends about $0.90$ kg less far from $m$ at each end. Percentage comparisons differ because the endpoint magnitudes differ.
\secInicio[][sec:q14]
}

\needspace{38\baselineskip}
\sectionScored{Production cost savings}{C9-\sml{2}}{sec:q15}
{
\textbf{C9 -- Inverse lower-tail cutoff.} A sample of $n=12$ batches has cost deviation from budget, in dollars, with mean $m=-1.6$ and approximate standard deviation $a=4$. Compare $X_T=-1.6+4T$, $T\sim t_{n-1}$, with $X_N=-1.6+4Z$.
\QuestionsTasks[sec:q15]{
{9}/{sol:q15}/{Find the cutoff $x$ below which $5\%$ of deviations fall. Find both model results, calculate $E_A$, $E_{Pt}$, and $E_{Pn}$. \score{2}}
}
\needspace{5\baselineskip}
\phantomsection
\label{sol:q15}
\subsection*{C9 -- Problem 15: Inverse lower-tail cutoff solution \score{2}}
Here $\nu=12-1=11$. Standardize the unknown cutoff(s) using $c=(x-(-1.6))/4$. The lower-tail target is $0.05$.
Row $11$, column $\alpha=0.05$, gives $t^*=1.796$. 
Obtain the negative cutoff by symmetry, giving $c=-1.796$.
For normal, row $1.6$, columns $0.04$ and $0.05$, gives $0.9495$ and $0.9505$. The negative cumulative entries are $0.0505$ and $0.0495$, equally close to $0.0500$. The tie rule selects $c_N=-1.64$.
Back-transform, retaining unrounded values:
\[
\begin{aligned}
x_T&=-1.6+4(-1.796)=-8.784\approx\boxed{-8.78\text{ dollars}}.
\end{aligned}
\]
\[
\begin{aligned}
x_N&=-1.6+4(-1.64)=-8.16\approx\boxed{-8.16\text{ dollars}}.
\end{aligned}
\]
Use unrounded contextual results:
\[
\begin{aligned}
E_{A}&=|-8.784-(-8.16)|=0.624\approx\boxed{0.62\text{ dollars}},\\
E_{Pt}&=\frac{100(0.624)}{|-8.784|}\approx\boxed{7.10\%},\\
E_{Pn}&=\frac{100(0.624)}{|-8.16|}\approx\boxed{7.65\%}.
\end{aligned}
\]
About $5\%$ of batches save more than $\$8.78$ under $t$, or $\$8.16$ under normal. The normal threshold is about $\$0.62$ closer to the mean deviation.
\secInicio[][sec:q15]
}

\needspace{38\baselineskip}
\sectionScored{Order processing cutoff}{C9-\sml{2}}{sec:q16}
{
\textbf{C9 -- Inverse upper-tail cutoff.} A sample of $n=21$ orders has processing-time deviation from target, in minutes, with mean $m=1.8$ and approximate standard deviation $a=5$. Compare $X_T=1.8+5T$, $T\sim t_{n-1}$, with $X_N=1.8+5Z$.
\QuestionsTasks[sec:q16]{
{9}/{sol:q16}/{Find the cutoff $x$ exceeded by $10\%$ of deviations. Find both model results, calculate $E_A$, $E_{Pt}$, and $E_{Pn}$. \score{2}}
}
\needspace{5\baselineskip}
\phantomsection
\label{sol:q16}
\subsection*{C9 -- Problem 16: Inverse upper-tail cutoff solution \score{2}}
Here $\nu=21-1=20$. Standardize the unknown cutoff(s) using $c=(x-(1.8))/5$. The right-tail target is $0.1$.
Row $20$, column $\alpha=0.1$, gives $t^*=1.325$. 
The positive cutoff is $c=1.325$.
Convert right tail $0.10$ to cumulative target $0.9000$. For normal, row $1.2$, column $0.08$, gives $\Phi(1.28)=0.8997$, closer to $0.9000$ than $\Phi(1.29)=0.9015$ (column $0.09$). Thus $c_N=1.28$.
Back-transform, retaining unrounded values:
\[
\begin{aligned}
x_T&=1.8+5(1.325)=8.425\approx\boxed{8.43\text{ minutes}}.
\end{aligned}
\]
\[
\begin{aligned}
x_N&=1.8+5(1.28)=8.2\approx\boxed{8.20\text{ minutes}}.
\end{aligned}
\]
Use unrounded contextual results:
\[
\begin{aligned}
E_{A}&=|8.425-(8.2)|=0.225\approx\boxed{0.23\text{ minutes}},\\
E_{Pt}&=\frac{100(0.225)}{|8.425|}\approx\boxed{2.67\%},\\
E_{Pn}&=\frac{100(0.225)}{|8.2|}\approx\boxed{2.74\%}.
\end{aligned}
\]
About $10\%$ of orders take more than $8.43$ minutes above target under $t$, or $8.20$ minutes under normal. The normal cutoff is about $0.23$ minutes lower.
\secInicio[][sec:q16]
}

\needspace{48\baselineskip}
\sectionScored{Research option}{C10-\sml{6}}{sec:q17}
{
\textbf{Exam:}
\QuestionsTasks[sec:q17]{
{10}/{sol:q17a}/{Select one model and describe plausible real-life data in 1--2 sentences. \score{6}}
}
\begin{enumerate}[label=\textbf{(\alph*)},leftmargin=*]
\item Linear probability density over a finite interval.
\item Exponential distribution (waiting times).
\item Normal distribution (measurements).
\end{enumerate}

\textbf{Online report:} Submit a PDF through \textbf{Open} using real economic or financial data and the selected model. Include simulation, sampling, portfolio risk, methods, assumptions, limitations, and evidence-based conclusions. The report is the C10 evidence.

\needspace{6\baselineskip}
\phantomsection
\label{sol:q17a}
\subsection*{C10 -- Acceptable alternative brief responses \score{6}}

This selection task has no unique answer. For this worked response, select \textbf{(b)} and use its two-sentence response below.
A student submits \textbf{only one} 1--2 sentence response; (a) and (c) are clearly identified \textbf{alternatives}:
\begin{enumerate}[label=\textbf{(\alph*)},leftmargin=*]
\item \textbf{Linear density.} I would study customers' spending in thousands of Colombian pesos on a promotional service priced continuously between 20 and 60 thousand pesos, proposing a linearly decreasing density over $[20,60]$. I would check whether real observations support that proposed model.
\item \textbf{Exponential waiting times.} I would study the minutes between payment requests reaching a small business during a stable trading period, proposing an exponential waiting-time model if arrivals are independent and occur at a constant average rate. I would check that assumption with real observations.
\item \textbf{Normal measurements.} I would study the fill volumes in millilitres of bottles sold by a beverage producer, proposing a normal model during stable production. I would check whether real measurements support that model.
\end{enumerate}
Each describes possible real-life economic data and a \emph{proposed} model requiring validation;
none claims that uncollected data follow the model. Other plausible responses are acceptable.
Each example is a brief exam response, not an empirical analysis or a completed online report.
The exam prints six points beside the brief selection and identifies the separate online report as the C10 evidence; no separate report point allocation is stated.

\textbf{Separate online PDF report (not the brief exam response).}
Submit through \textbf{Open}, focusing on the \textbf{one selected option}.
Investigate an economic or financial question using real data and the selected continuous model.
Include simulation, sampling analysis, and portfolio-risk measures; justify methods and present results,
assumptions, limitations, and evidence-based conclusions. These are the broader C10 requirements;
no completed report, dataset, simulation, portfolio analysis, or empirical conclusion is supplied by these examples.

\secInicio[][sec:q17]
}

\newpage
\section{Summarized Solutions}
\begingroup
\small
\setlength{\parskip}{5pt}
\setlength{\tabcolsep}{5pt}
\renewcommand{\arraystretch}{1.22}
\textbf{Notation and conventions.}
$\nu=n-1$, $c=(x-m)/a$, $x=m+ac$;
$X_T=m+aT$, $T\sim t_\nu$; $X_N=m+aZ$, $Z\sim N(0,1)$.
Use the stated approximate SD directly as scale $a$: the models share a scale,
not an actual variance. These are signed-deviation models; do not use
$a/\sqrt n$ or interpret inverse intervals as confidence intervals for a population mean.
Negative contextual values mean below the relevant reference; positive values mean above it.

All results are approximations from the supplied printed tables. No interpolation,
CDF/inverse-CDF calculators or exact distribution values are used.
Direct normal bounds are rounded to two decimals before lookup; inverse normal
cutoffs use the closest printed CDF entry, with the $0.0500/0.9500$ tie selecting
$|z|=1.64$. Retain three-decimal $t$ and two-decimal $z$ cutoffs.
In tables marked $\approx$, numerical results are table-based/rounded;
$=$ denotes identities or arithmetic with selected working values.

\subsection*{C8}
\textbf{A. Direct probabilities --- Problems 1--5}

\begin{tabularx}{\linewidth}{@{}lr>{\raggedright\arraybackslash}Xllr@{}}
\toprule
Problem & $\nu$ & Event & Bound(s) & Area calculation & $P$ ($\approx$)\\
\midrule
1 & 9 & $X<-7.048\text{ or }X>7.532$ & \shortstack[l]{$-2.262$\\$1.383$} & $0.025+0.1$ & $0.125$\\
2 & 20 & $6.198<X<10.95$ & \shortstack[l]{$0.533$\\$1.325$} & $0.3-0.1$ & $0.2$\\
3 & 12 & $X<-4.624$ & \shortstack[l]{$-1.356$} & $0.1$ & $0.1$\\
4 & 18 & $X>-8.272$ & \shortstack[l]{$-0.534$} & $1-0.3$ & $0.7$\\
5 & 11 & $-6.188<X<-2.42$ & \shortstack[l]{$-1.796$\\$-0.540$} & $0.3-0.05$ & $0.25$\\
\bottomrule
\end{tabularx}
Event units (1--5): thousands of dollars, kWh, minutes, dollars, tonnes.

\medskip
\textbf{B. Inverse cutoffs and intervals --- Problems 6--10}

\begin{tabularx}{\linewidth}{@{}lr>{\raggedright\arraybackslash}Xlll@{}}
\toprule
Problem & $\nu$ & Tail(s) & Signed $t$ & Result ($\approx$) & Units\\
\midrule
6 & 19 & left 0.1 & \shortstack[l]{$-1.328$} & \shortstack[l]{$-11.50$} & dollars\\
7 & 14 & left 0.025; right 0.025 & \shortstack[l]{$-2.145$\\$2.145$} & \shortstack[l]{$-4.56$\\$3.16$} & g\\
8 & 22 & left 0.01; right 0.1 & \shortstack[l]{$-2.508$\\$1.321$} & \shortstack[l]{$-16.85$\\$6.13$} & minutes\\
9 & 25 & right 0.025 & \shortstack[l]{$2.060$} & \shortstack[l]{$17.42$} & minutes\\
10 & 16 & left 0.05; right 0.05 & \shortstack[l]{$-1.746$\\$1.746$} & \shortstack[l]{$-11.71$\\$19.71$} & tonnes\\
\bottomrule
\end{tabularx}
For pairs, results are $(L,U)$ in order. Problem 8: outside $=0.01+0.10=0.11$; inside $=0.89$.
\newpage
\subsection*{C9}
\textbf{Comparison measures.}
$E_A=|r_T-r_N|$, $E_{Pt}=100E_A/|r_T|$,
$E_{Pn}=100E_A/|r_N|$.
For direct problems, $r_T=p_T$, $r_N=p_N$; for inverse problems, $r_T,r_N$
are contextual cutoffs after back-transformation. These compare model results,
not errors against observed truth. $E_{Pt},E_{Pn}$ are relative percentages;
a probability difference of $E_A$ corresponds to $100E_A$ percentage points.

\textbf{A. Direct model comparisons --- Problems 11--13}

\begin{tabularx}{\linewidth}{@{}Xrrrrrr@{}}
\toprule
Problem & $\nu$ & $p_T$ & $p_N$ & $E_A$ & $E_{Pt}$ & $E_{Pn}$\\
\midrule
11 & 16 & 0.0250 & 0.0170 & 0.0080 & 32.00\% & 47.06\%\\
12 & 18 & 0.7500 & 0.7633 & 0.0133 & 1.77\% & 1.74\%\\
13 & 15 & 0.7000 & 0.7054 & 0.0054 & 0.77\% & 0.77\%\\
\bottomrule
\end{tabularx}
All numerical results $\approx$; probabilities and $E_A$ are unitless.

Problem 11: $c=-2.120\to z=-2.12$.\par
Problem 12: $c=-0.862\to z=-0.86$, $c=1.734\to z=1.73$.\par
Problem 13: $c=-0.536\to z=-0.54$.\par
\medskip
\textbf{B. Inverse model comparisons --- Problems 14--16}

\begin{tabularx}{\linewidth}{@{}Xrrrrrl@{}}
\toprule
Problem/endpoint & $r_T$ & $r_N$ & $E_A$ & $E_{Pt}$ & $E_{Pn}$ & Units\\
\midrule
14 / $L$ & -7.52 & -6.62 & 0.90 & 11.97\% & 13.60\% & kg\\
14 / $U$ & 11.92 & 11.02 & 0.90 & 7.55\% & 8.17\% & kg\\
15 / $x$ & -8.78 & -8.16 & 0.62 & 7.10\% & 7.65\% & dollars\\
16 / $x$ & 8.43 & 8.20 & 0.23 & 2.67\% & 2.74\% & minutes\\
\bottomrule
\end{tabularx}
All numerical results $\approx$; contextual cutoffs and $E_A$ have the stated units.

Problem 14 ($\nu=13$): $c_T=-2.160,\ z=-1.96$, $c_T=2.160,\ z=1.96$.\par
Problem 15 ($\nu=11$): $c_T=-1.796,\ z=-1.64$.\par
Problem 16 ($\nu=20$): $c_T=1.325,\ z=1.28$.\par
Use unrounded back-transformed values for comparisons. For Problem 16, $r_T=8.425$, $r_N=8.20$, $E_A=0.225$ in selected working values.
\needspace{23\baselineskip}
\subsection*{C10}
\par\smallskip
\begin{tabularx}{\linewidth}{@{}p{3.1cm} >{\raggedright\arraybackslash}X@{}}
\toprule
Component & Summarized response or requirement \\
\midrule
Selected option & \textbf{(b), exponential waiting times.} \\
Brief exam response & I would study the minutes between payment requests reaching a small business during a stable trading period, proposing an exponential waiting-time model if arrivals are independent and occur at a constant average rate. I would check that assumption with real observations. \\
Alternatives & \textbf{(a)}, linear density over a finite interval; \textbf{(c)}, normal measurements. Acceptable alternatives; students submit \textbf{only one 1--2 sentence response}. Other plausible responses are acceptable. \\
\midrule
\shortstack[l]{Separate online\\report} & \textbf{PDF submitted through Open}, using real economic or financial data and the \textbf{one selected model}. Include simulation, sampling analysis, portfolio-risk measures, justified methods, results, assumptions, limitations, and evidence-based conclusions. \\
\bottomrule
\end{tabularx}
\par\smallskip
The brief response proposes data and a model to validate; it supplies no completed
report or empirical findings. The online report is the C10 evidence;
no separate report point allocation is stated.
\par
\endgroup

\end{document}
